\chapter{System Design}
\label{chap:system_design}

This chapter explores the high-level design of the benchmarking suite. 
Our system is based on a simulated exchange designed to process trades from 
multiple concurrent clients. 
To understand the system from first principles, we have engineered a 
conventional exchange pipeline mostly from scratch. 
This multithreaded architecture parses FIX messages, passes them to a custom 
lock-free queue for thread hand-off, and finally dispatches them to the 
matching engine. 

While every component in this pipeline contributes to total system latency, 
the order book is typically the primary stateful core of a 
matching engine~\citep{bilokon2023cpp}. As the heart of the exchange, 
the order book must maintain price-time priority while processing 
high-frequency updates, making the choice of data structure the most 
significant factor for engine scalability~\citep{harrisequity}. For this 
reason, this report focuses on exploring the performance trade-offs of 
varied order book layouts. 

When designing a benchmarking suite, it is essential that all measurements 
are scientific and fair. We maintain this rigour by validating that all 
components operate with semantic parity, ensuring that latency 
differences are caused by architectural efficiency rather than logical 
shortcuts.

\section{Formal Requirements}
\label{sec:system_formal_requirements}
To ensure academic rigour and verify that the system performs as intended, 
we use a formal requirements framework. This approach is inspired by 
traditional software engineering methodologies for requirements engineering 
\citep{sommerville2015} and the rigorous systems benchmarking principles 
advocated by \citet{gregg2013systems}. These requirements serve as 
the analytical foundation for the project and provide the objective 
success criteria for the performance evaluation in Chapter 5. By 
establishing these benchmarks early, we ensure that the subsequent 
analysis of order book implementations is grounded in measurable and 
verifiable standards.


\subsection{Functional Requirements}
\begin{table}[H]
    \centering
    \small
    \begin{tabular}{lp{5cm}p{6.5cm}}
        \toprule
        \textbf{ID} & \textbf{Requirement} & \textbf{Testable Criterion} \\
        \midrule
        \phantomsection\label{fr:1} FR-1 & Implement five distinct order book data structures & Five implementations present and benchmarkable \\
        \phantomsection\label{fr:2} FR-2 & Support direct-mode benchmarking isolating engine latency & Direct mode runs without network stack involvement \\
        \phantomsection\label{fr:3} FR-3 & Support gateway-mode benchmarking with end-to-end FIX protocol overhead & Gateway mode includes TCP, FIX parsing, and queue handoff \\
        \phantomsection\label{fr:4} FR-4 & Support MPSC contention benchmarking with configurable producer counts & MPSC mode scales across producer configurations \\
        \phantomsection\label{fr:5} FR-5 & Emit structured CSV output for all benchmark runs & CSV files produced with consistent schema per run \\
        \phantomsection\label{fr:6} FR-6 & Include correctness test suite verifying semantic parity across implementations & All implementations produce identical outputs for equivalent inputs \\
        \bottomrule
    \end{tabular}
    \caption{Functional Requirements (FR)}
    \label{tab:functional_requirements}
\end{table}

\subsection{Non-Functional Requirements}
\begin{table}[H]
    \centering
    \small
    \begin{tabular}{lp{5cm}p{6.5cm}}
        \toprule
        \textbf{ID} & \textbf{Requirement} & \textbf{Testable Criterion} \\
        \midrule
        \phantomsection\label{nfr:1} NFR-1 & Correctness: semantic parity across all implementations & No variant has special-case logic biasing results \\
        \phantomsection\label{nfr:2} NFR-2 & Fairness: identical workload parameters across all implementations & Same run count, order count, scenario selection for all books \\
        \phantomsection\label{nfr:3} NFR-3 & Isolation: direct mode measures only algorithmic/data-structure latency & No transport overhead in direct-mode code path \\
        \phantomsection\label{nfr:4} NFR-4 & Realism: gateway mode includes network, queue, parsing, and engine latency & End-to-end timing supports dissertation's masking argument \\
        \phantomsection\label{nfr:5} NFR-5 & Reproducibility: pipeline regenerates canonical dataset without undocumented steps & Fresh run produces consistent artifact structure \\
        \phantomsection\label{nfr:6} NFR-6 & Statistical rigour: mean, standard deviation, and tail metrics reported & Variance metrics present and consistently derived \\
        \phantomsection\label{nfr:7} NFR-7 & Performance isolation: CPU pinning and scheduler isolation applied & Pinning policy documented and verified in results \\
        \phantomsection\label{nfr:8} NFR-8 & Portability discipline: absolute latency claims bounded to WSL2/AMD environment & Relative comparisons defensible; absolute numbers environment-scoped \\
        \phantomsection\label{nfr:9} NFR-9 & Claim discipline: no measurement overstated beyond evidence & Cache and allocator causality stated as correlation, not proven cause \\
        \bottomrule
    \end{tabular}
    \caption{Non-Functional Requirements (NFR)}
    \label{tab:non_functional_requirements}
\end{table}

\section{Design Goals and Constraints}
The primary design goal of this benchmarking suite is to minimise 
measurement overhead and avoid the `observer effect' (or `Heisenberg 
effect'), where the instrumentation itself alters the system's performance 
characteristics~\citep{gregg2013systems}. Ensuring that the act of 
measuring latency does not significantly inflate the results is essential 
for maintaining the scientific validity of the data.

Additionally, the project must navigate the technical constraints of the 
Linux/WSL2 environment. While this virtualised setup introduces 
additional jitter and overhead compared to bare-metal hardware, these 
factors are explicitly documented and accounted for throughout the 
evaluation. By acknowledging these boundaries, we ensure that our 
conclusions are grounded in the specific realities of the test 
environment.





\section{High-Level Architecture}
The exchange is designed as a modular system~\citep{parnas1972} using the `Ports and 
Adapters' (Hexagonal) architecture~\citep{velepucha2023hexagonal,cockburn_hexagonal}. This 
approach isolates the core trading logic from external infrastructure, 
ensuring that the system's performance characteristics remain 
deterministic and decoupled from transport or logging overheads.

\subsection{Architectural Pattern: Ports and Adapters}
In this design, the core \textit{Matching Engine} (\texttt{src/core}) 
functions as the central "hexagon" containing the domain logic. It remains 
unaware of whether it is communicating with a real network client or a 
simulated benchmark feed. Communication with the engine happens 
through \textit{Ports} - abstract interfaces that define what 
the engine requires from the outside world.

The most critical port in this system is the \texttt{IOrderBook} interface. 
By defining a stable API for order insertion, cancellation, and matching, 
the engine can interact with any data structure implementation (the 
\textit{Adapters}) without modification. This encapsulation is the 
foundation of our scientific fairness (\NFR{2}); it allows us to swap 
an \texttt{ArrayOrderBook} for a \texttt{MapOrderBook} while keeping the 
surrounding benchmark harness and engine logic identical.

\subsection{System Modules}
As illustrated in Figure~\ref{fig:hexagonal_architecture}, the system 
is composed of four primary modules:
\begin{itemize}
    \item \textbf{Order Generator}: A driving adapter that generates 
    synthetic market workloads for the system.
    \item \textbf{Gateway}: A network-facing adapter that parses FIX 
    messages and manages the transition between network and engine threads.
    \item \textbf{Matching Engine}: The core domain logic that processes 
    the central limit order book.
    \item \textbf{Telemetry}: A driven adapter responsible for capturing 
    nanosecond-resolution metrics at each stage of the pipeline.
\end{itemize}

\begin{figure}[H]
    \centering
    \resizebox{0.92\textwidth}{!}{\begin{tikzpicture}[
    node distance=2cm, 
    auto,
    block/.style={rectangle, draw, fill=blue!10, text width=7.5em, text centered, rounded corners, minimum height=3em},
    line/.style={draw, -latex'},
    interface/.style={draw, diamond, fill=green!10, text width=6.5em, text centered, inner sep=0pt}
]

    % Nodes
    \node [block] (gen) {Order Generator};
    \node [block, below of=gen, xshift=-3cm] (direct) {Direct Benchmark};
    \node [block, below of=gen, xshift=3cm] (client) {Mock Client (Gateway)};
    \node [block, below=1.5cm of client] (gateway) {TCP Gateway (Server)};
    \node [interface, below of=gateway, xshift=-3cm] (iface) {\texttt{IOrderBook}};
    \node [block, below of=iface] (book) {Implementation (e.g. Array)};
    \node [block, right=2cm of book] (telemetry) {Metric Collection};

    % Paths
    \path [line] (gen) -- (direct);
    \path [line] (gen) -- (client);
    \path [line] (client) -- node[font=\scriptsize]{FIX/TCP} (gateway);
    \path [line] (direct) |- (iface);
    \path [line] (gateway) -- (iface);
    \path [line] (iface) -- (book);
    \path [line] (book) -- (telemetry);
    \path [line] (gateway) -| (telemetry);

\end{tikzpicture}
}
    \caption{Data-Flow Overview: Pipeline stages from generation to telemetry.}
    \label{fig:high_level_architecture}
\end{figure}

\clearpage
\begin{figure}[H]
    \centering
    \resizebox{0.92\textwidth}{!}{\begin{tikzpicture}[
    node distance=3.5cm,
    block/.style={rectangle, draw, fill=blue!10, text width=7em, text centered, rounded corners, minimum height=3em, font=\scriptsize},
    port/.style={draw, fill=green!10, minimum width=2.5cm, font=\bfseries\scriptsize},
    hex/.style={draw, regular polygon, regular polygon sides=6, minimum size=6cm, fill=gray!2, thick}
]

    % The Hexagon (Core Domain)
    \node [hex] (hexagon) at (0,0) {};
    \node [font=\bfseries] at (0,1.2) {Matching Engine};
    \node [font=\scriptsize, text width=3.5cm, align=center] at (0,0.6) {Core Trading Domain};
    \node [font=\tiny\itshape, text width=3.5cm, align=center, gray] at (0,0.2) {(\texttt{src/core})};

    % Ports (Positioned on the edges of the hexagon logic)
    \node [port, fill=green!15] (input_port) at (-2.6, -0.2) {Driving Port};
    \node [port, fill=yellow!15] (output_port) at (2.6, -0.2) {Driven Port};
    \node [port, fill=orange!15] (lob_port) at (0, -2.4) {IOrderBook Port};

    % Adapters (Outside Hexagon - clear separation)
    
    % Driving Adapters (Inputs)
    \node [block, left=1.2cm of input_port, yshift=1cm] (gen) {\textbf{Order Generator}\\(Scenario Feed)};
    \node [block, left=1.2cm of input_port, yshift=-1cm] (gw) {\textbf{TCP Gateway}\\(FIX/TCP)};

    % Driven Adapters (Outputs)
    \node [block, right=1.2cm of output_port, yshift=1cm] (tele) {\textbf{Telemetry Suite}\\(RDTSC Timestamps)};
    \node [block, right=1.2cm of output_port, yshift=-1cm] (csv) {\textbf{CSV Exporter}\\(Metrics Persistence)};

    % The Interchangeable Component
    \node [block, below=1cm of lob_port] (lob_impl) {\textbf{Data Structure Adapter}\\(Array, Map, Pool, etc.)};

    % Connecting Adapters to Ports
    \draw [thick, <->] (gen) -- (input_port);
    \draw [thick, <->] (gw) -- (input_port);
    \draw [thick, <->] (tele) -- (output_port);
    \draw [thick, <->] (csv) -- (output_port);
    \draw [thick, ->] (lob_port) -- (lob_impl);

    % Boundary Labels
    \node [font=\tiny\itshape, gray] at (-4, 2.5) {External Actors};
    \node [font=\tiny\itshape, gray] at (4, 2.5) {Infrastructure};

\end{tikzpicture}
}
    \caption{Hexagonal Architecture (Ports and Adapters): Clear separation between the matching domain and interchangeable adapters.}
    \label{fig:hexagonal_architecture}
\end{figure}

\section{The IOrderBook Interface and Polymorphism}
The core of our modular design is the \texttt{IOrderBook} interface, 
implemented as a C++ abstract base class (\texttt{src/core/i\_order\_book.hpp}).
This interface defines the contract that any order book data 
structure must satisfy to be compatible with the matching engine. It 
standardises essential operations such as order insertion 
(\texttt{addOrder}), cancellation (\texttt{cancelOrder}), and the matching 
logic (\texttt{match}).

Using a polymorphic interface ensures that each implementation is treated 
as a `black box' by the benchmark harness. This is critical for 
scientific fairness (\NFR{1}); it prevents the benchmarking logic from 
using implementation-specific shortcuts and ensures that every data structure 
is subjected to the exact same call overhead (i.e. the nanosecond-scale 
latency of virtual method table lookups). Separating the engine's logic 
from the book's implementation ensures that our results only reflect 
the efficiency of the data structures being tested.

\section{Benchmarking Design: The Two-Mode Philosophy}
To validate our hypotheses scientifically, we adopt a dual-mode 
benchmarking philosophy that separates algorithmic efficiency from 
overhead. This design is necessary so that our measurements capture 
both the theoretical potential of our data structures and their relative 
value in a production-like environment.

The \textbf{Direct Mode} serves as a controlled micro-benchmark. It 
executes the matching engine logic synchronously within the benchmark 
thread, bypassing all networking and asynchronous layers. This isolation is 
essential for uncovering the `pure' performance characteristics of the 
underlying order book, allowing us to see latency differences that are 
usually hidden by the surrounding system.

The \textbf{Gateway Mode} provides realism by including the 
full TCP/IP stack (via loopback), FIX protocol parsing, and the 
asynchronous hand-off between network and engine threads. 
Measuring the system in this mode is critical for exploring the \textit{Masking 
Effect} proposed in Hypothesis~\ref{hyp:h3}. 
As we will evaluate in Chapter 5, the constant-time overhead of these 
networking layers obscures the relative performance advantages of faster 
data structures. 
By designing our benchmark suite with both modes, we can differentiate between 
a structure's peak efficiency and its practical impact on end-to-end latency.

\section{Concurrency and Threading Strategy}
To maintain low latency, the system must separate high-latency I/O 
operations - such as network packet polling, TCP/IP stack traversal, and 
FIX message parsing - from the core matching engine. 
This multi-threaded architecture prevents slower network components from 
blocking the execution of trades, which is a standard pattern for 
high-performance sequence processing in industrial matching 
engines~\citep{thompson2011disruptor}.

At the core of the concurrency strategy is the \textbf{Producer-Consumer} 
hand-off. The Gateway thread (Producer) handles the I/O operations 
described above and pushes parsed orders into a lock-free queue, while the 
Matching Engine thread (Consumer) pulls and executes them. For the 
standard gateway mode, we employ a \textbf{Single-Producer Single-Consumer 
(SPSC)} queue, which provides wait-free properties for the engine. In 
scenarios with multiple concurrent clients, the system can swap this for a 
Multi-Producer Multi-Consumer (MPMC) adapter from the \texttt{moodycamel} 
library.

To ensure deterministic performance, we use \textbf{CPU Pinning} (Thread 
Affinity). 
By pinning the engine thread to a specific physical core, we prevent the 
operating system from migrating the process across different CPUs. 
This elimination of context-switching jitter and cache-migration overhead is 
vital for mitigating the \textit{Observer Effect}. 
It ensures that our measurements represent 'bare-metal' performance as closely 
as possible within a Linux environment.

\section{Implementation Parameters}
\label{sec:implementation_parameters}
To ensure reproducibility (\NFR{5}), we specify the following design 
parameters used across all benchmarked variants: 

\begin{itemize}
    \item \textbf{Array Layout}: 100 levels per side (1.6~KB total). This size is intentionally chosen to fit entirely within the 32~KB L1 Data Cache of the test hardware.
    \item \textbf{Pool Capacity}: 1 million pre-allocated \texttt{Order} structures (approx. 44~MB). This volume intentionally exceeds the 32~MB L3 cache to test the performance of slab recycling after cache exhaustion.
    \item \textbf{Hybrid Hot-Path}: 20-level threshold for vector storage. This represents approximately 80\% of typical market activity according to our scenario generation profiles.
\end{itemize}
